\documentclass[10pt,a4paper]{amsart}
\usepackage[margin=19mm]{geometry}
\usepackage{amsmath,amssymb,mathtools}
\usepackage[scr=rsfs,cal=euler]{mathalfa}
\usepackage{xfrac}
\usepackage[unicode,hidelinks,bookmarks=false]{hyperref}
\newcommand{\ie}{i.e.\ }
\newcommand{\cf}{cf.\ }
\newcommand{\resp}{respectively\ }
\newcommand{\Subsection}{\subsection}
\providecommand{\nobreakdash}{\nobreak\hbox{-}}
\setlength{\emergencystretch}{2em}
\multlinegap0pt
\usepackage{amssymb}
\usepackage{algorithm}\usepackage{algpseudocode}
\usepackage{tikz}
\usepackage{multirow}
\usepackage{booktabs}
\usepackage{bm}
\newtheorem{lem}{Lemma}
\title{Phase Transitions, Non-Extremality (Reconstruction), and Markov Entropy Rate for the Mixed Spin-$(s,\tfrac12)$ Ising Model on a Cayley Tree of Order Three}
\author{Hasan Akin}
\date{Source: arXiv:2602.13412v1 (CC BY 4.0)}
\begin{document}
\maketitle

\noindent\emph{Source and licence.} Phase Transitions, Non-Extremality (Reconstruction), and Markov Entropy Rate for the Mixed Spin-$(s,\tfrac12)$ Ising Model on a Cayley Tree of Order Three. Version arXiv:2602.13412v1 (CC BY 4.0), distributed by arXiv under the Creative Commons Attribution 4.0 International licence, http://creativecommons.org/licenses/by/4.0/, as recorded in the arXiv licence metadata for this exact version and shown on the abstract page https://arxiv.org/abs/2602.13412v1. Full licence notice: \texttt{licenses/arxiv-2602.13412-CC-BY-4.0.txt}.\par\smallskip

\noindent\textit{Editorial scope.} Counted unit: the article's lem lem:Dobrushin-general (source0.tex lines 1363--1451). The article's complete original proof of it is delivered with it (source0.tex lines 1453--1547, one \texttt{proof} environment of 95 lines). No mathematics is written, repaired or re-derived: the statement and the proof are the article's own lines, in the article's own notation, and no retained line is substituted or re-flowed.  Retained uncounted as the source-local context the proof needs: lines 1122--1136; lines 1208--1220; lines 1338--1354.  Nothing was omitted from any retained span, so the package carries no omission record.  No retained line contains a citation, so the wrapper carries no bibliography and no bibliography entry.  No retained line contains a figure environment, an image-inclusion command or drawing code, so no image is delivered and the package declares no asset dependency.  Editorial formatting only, in addition: an AMS \texttt{amsart} preamble that restates the article's own declarations and macros used by the retained lines, namely the environments \texttt{lem} on separate counters, together with the standard packages those lines need.  The retained environments print in this document as Lemma 1 in order of appearance, whereas the article prints the same items with the numbering of its own full document.  The complete native source of the article as distributed in the arXiv e-print for this exact version is bundled unchanged as \texttt{sources/194633/source0.tex}.
\begin{align}\label{eq:StocM-s=5-P}
\mathbb{P}^{(5)}(\varphi)
=
\begin{pmatrix}
\dfrac{\varphi^{10}}{\varphi^{10}+1} & \dfrac{1}{\varphi^{10}+1}\\[4pt]
\dfrac{\varphi^{8}}{\varphi^{8}+1} & \dfrac{1}{\varphi^{8}+1}\\[2pt]
\vdots & \vdots\\[2pt]
\dfrac{\varphi^{2}}{\varphi^{2}+1} & \dfrac{1}{\varphi^{2}+1}\\[4pt]
\dfrac{1}{2} & \dfrac{1}{2}\\[4pt]
\dfrac{1}{1+\varphi^{2}} & \dfrac{\varphi^{2}}{1+\varphi^{2}}\\[2pt]
\vdots & \vdots\\[2pt]
\dfrac{1}{1+\varphi^{8}} & \dfrac{\varphi^{8}}{1+\varphi^{8}}\\[4pt]
\dfrac{1}{1+\varphi^{10}} & \dfrac{\varphi^{10}}{1+\varphi^{10}}
\end{pmatrix}.  
\end{align}

\begin{align}\label{eq:Stochas-Mat-Q(5)}
\mathbb{Q}^{(5)}(\varphi)
=
\frac{1}{S(\varphi)}
\left(
\begin{array}{ccccccccccc}
h_5(\varphi) & h_4(\varphi) & h_3(\varphi) & h_2(\varphi) & h_1(\varphi) & 1
             & g_1(\varphi) & g_2(\varphi) & g_3(\varphi) & g_4(\varphi) & g_5(\varphi)\\
g_5(\varphi) & g_4(\varphi) & g_3(\varphi) & g_2(\varphi) & g_1(\varphi) & 1
             & h_1(\varphi) & h_2(\varphi) & h_3(\varphi) & h_4(\varphi) & h_5(\varphi)
\end{array}
\right),
\end{align}

\begin{align}\label{eqs:Dobrushin-def1a}
\tau_{\mathbb{P}^{(s)}}(\varphi)
=
\frac{1}{2}\max_{i,j}
\left\{
\sum_{\ell=1}^{2}
\bigl|P_{i\ell} - P_{j\ell}\bigr|
\right\},
\qquad
\tau_{\mathbb{Q}^{(s)}}(\varphi)
=
\frac{1}{2}\max_{i,j}
\left\{
\sum_{\ell=1}^{2s+1}
\bigl|Q_{i\ell} - Q_{j\ell}\bigr|
\right\}.  
\end{align}

\begin{lem}\label{lem:Dobrushin-general}
Let $s\in\{1,2,3,4,5\}$ and let $\tau_{\mathbb{P}^{(s)}}(\varphi)$ and 
$\tau_{\mathbb{Q}^{(s)}}(\varphi)$ denote the Dobrushin coefficients of the
stochastic matrices $\mathbb{P}^{(s)}(\varphi)$ and $\mathbb{Q}^{(s)}(\varphi)$
corresponding to the mixed spin $(s,\tfrac12)$ model. Then, for all
$\varphi>0$, from \eqref{eqs:Dobrushin-def1a}, one has
\begin{equation}\label{eq:tauP-s-phi-def}
\tau_{\mathbb{P}^{(s)}}(\varphi)
=
\frac{\bigl|\varphi^{2s} - 1\bigr|}{\varphi^{2s} + 1},
\end{equation}
and
\begin{equation}\label{eq:tauQ-s-phi-def}
\tau_{\mathbb{Q}^{(s)}}(\varphi)
=
\frac{1}{8\,S_s(\varphi)}
\sum_{n=1}^{s}
\frac{\bigl(\varphi^{2n}+1\bigr)^3\,\bigl|\varphi^{2n}-1\bigr|}{\varphi^{4n}},
\qquad
S_s(\varphi)
=
1+\frac18\sum_{n=1}^{s}\frac{\bigl(\varphi^{2n}+1\bigr)^4}{\varphi^{4n}}.
\end{equation}
In particular, for $s=5$ the explicit expressions reduce to
\begin{equation}\label{eq:Dob-Coff-P(5)-a}
\tau_{\mathbb{P}^{(5)}}(\varphi)
= \frac{|\varphi^{10} - 1|}{\varphi^{10} + 1},
\end{equation}
and
\begin{equation}\label{eq:Dob-Coff-Q(5)-a}
\tau_{\mathbb{Q}^{(5)}}(\varphi)
=
\frac{1}{8\,S(\varphi)}
\sum_{n=1}^{5}
\frac{(\varphi^{2n}+1)^3\,\bigl|\varphi^{2n}-1\bigr|}{\varphi^{4n}},
\qquad
S_5(\varphi)
=
1+\frac18\sum_{n=1}^{5}\frac{(\varphi^{2n}+1)^4}{\varphi^{4n}}.
\end{equation} 
Moreover, the following properties hold.
\begin{enumerate}
  \item[(i)] \textbf{Symmetry.} For all $\varphi>0$,
  \begin{equation}\label{eq:tau-symmetry}
  \tau_{\mathbb{P}^{(s)}}(\varphi)
  =
  \tau_{\mathbb{P}^{(s)}}\!\left(\frac{1}{\varphi}\right),
  \qquad
  \tau_{\mathbb{Q}^{(s)}}(\varphi)
  =
  \tau_{\mathbb{Q}^{(s)}}\!\left(\frac{1}{\varphi}\right).
  \end{equation}
  
  \item[(ii)] \textbf{Bounds and values at $\varphi=1$.} For all $\varphi>0$,
  \begin{equation}\label{eq:tau-bounds}
  0 \,\le\, \tau_{\mathbb{P}^{(s)}}(\varphi) < 1,
  \qquad
  0 \,\le\, \tau_{\mathbb{Q}^{(s)}}(\varphi) < 1,
  \end{equation}
  and
  \begin{equation}\label{eq:tau-at-1}
  \tau_{\mathbb{P}^{(s)}}(1) = 0,
  \qquad
  \tau_{\mathbb{Q}^{(s)}}(1) = 0.
  \end{equation}
  
  \item[(iii)] \textbf{Continuity and asymptotics.} The functions
  $\varphi\mapsto\tau_{\mathbb{P}^{(s)}}(\varphi)$ and 
  $\varphi\mapsto\tau_{\mathbb{Q}^{(s)}}(\varphi)$ are continuous on $(0,\infty)$.
  Furthermore,
  \begin{equation}\label{eq:tau-asymptotics}
  \lim_{\varphi\to 0} \tau_{\mathbb{P}^{(s)}}(\varphi)
  =
  \lim_{\varphi\to\infty} \tau_{\mathbb{P}^{(s)}}(\varphi)
  = 1,
  \end{equation}
  and
  \begin{equation}\label{eq:tauQ-asymptotics}
  \lim_{\varphi\to 0} \tau_{\mathbb{Q}^{(s)}}(\varphi)
  =
  \lim_{\varphi\to\infty} \tau_{\mathbb{Q}^{(s)}}(\varphi)
  = 1.
  \end{equation}
\end{enumerate}
In particular, for each $s\in\{1,2,3,4,5\}$,
$\tau_{\mathbb{P}^{(s)}}(\varphi)$ and $\tau_{\mathbb{Q}^{(s)}}(\varphi)$ are
symmetric with respect to the point $\varphi=1$, vanish at $\varphi=1$, and
approach $1$ as $\varphi\to 0$ or $\varphi\to\infty$.
\end{lem}

\begin{proof}
We first treat the case $s=5$, for which the matrices
$\mathbb{P}^{(5)}(\varphi)$ and $\mathbb{Q}^{(5)}(\varphi)$ are given explicitly
in \eqref{eq:StocM-s=5-P} and \eqref{eq:Stochas-Mat-Q(5)}.

For $\mathbb{P}^{(5)}(\varphi)$ each row has the form $(a_i,1-a_i)$, so that
\[
\tau_{\mathbb{P}^{(5)}}(\varphi)
= \frac{1}{2}\max_{i,j} \left\{\sum_{\ell=1}^{2} |P_{i\ell} - P_{j\ell}|\right\}
= \max_{i,j} |a_i - a_j|,
\]
where $a_i = P_{i1}=\frac{1}{1+\varphi^{2i}}$ for $i\in\{-5,-4,\dots,5\}$.

Since $a_i$ is strictly decreasing in $i$ (for $\varphi>0$), the extreme values
occur at $i=-5$ and $i=5$, namely
\[
a_{\max} = a_{-5}=\frac{\varphi^{10}}{\varphi^{10} + 1},
\qquad
a_{\min} = a_{5}=\frac{1}{\varphi^{10} + 1}.
\]
Hence
\[
\max_{i,j} |a_i - a_j|
=
\left|\frac{\varphi^{10}}{\varphi^{10} + 1}
      -\frac{1}{\varphi^{10} + 1}\right|
=
\frac{|\varphi^{10} - 1|}{\varphi^{10} + 1},
\]
which yields \eqref{eq:Dob-Coff-P(5)-a}.

For the matrix $\mathbb{Q}^{(5)}(\varphi)$ in \eqref{eq:Stochas-Mat-Q(5)}, the
Dobrushin coefficient is
\[
\tau_{\mathbb{Q}^{(5)}}(\varphi)
=
\frac{1}{2}\max_{i,j}\left\{\sum_{\ell=1}^{11}
\bigl|\mathbb{Q}^{(5)}_{i\ell}(\varphi) - \mathbb{Q}^{(5)}_{j\ell}(\varphi)\bigr|\right\}.
\]
Since $\mathbb{Q}^{(5)}(\varphi)$ has only two rows, the maximum over $(i,j)$
reduces to this pair, and we obtain
\begin{equation}\label{eq:tauQ5-general}
\tau_{\mathbb{Q}^{(5)}}(\varphi)
=
\frac{1}{2}\sum_{\ell=1}^{11}
\bigl|\mathbb{Q}^{(5)}_{1\ell}(\varphi) - \mathbb{Q}^{(5)}_{2\ell}(\varphi)\bigr|.
\end{equation}

From \eqref{eq:Stochas-Mat-Q(5)}, we see that the differences between the two rows come in pairs
$h_n-g_n$ and $g_n-h_n$, whereas the middle column contributes zero. Hence
\eqref{eq:tauQ5-general} simplifies to
\[
\tau_{\mathbb{Q}^{(5)}}(\varphi)
=
\frac{1}{S(\varphi)}
\sum_{n=1}^{5} \bigl|h_n(\varphi) - g_n(\varphi)\bigr|.
\]
A direct calculation gives
\[
h_n(\varphi)-g_n(\varphi)
=
\frac{(\varphi^{2n}+1)^3}{8\varphi^{4n}}\bigl(\varphi^{2n}-1\bigr),
\]
so that
\[
\bigl|h_n(\varphi)-g_n(\varphi)\bigr|
=
\frac{(\varphi^{2n}+1)^3}{8\varphi^{4n}}\bigl|\varphi^{2n}-1\bigr|.
\]
Substituting this into the previous expression yields
\[
\tau_{\mathbb{Q}^{(5)}}(\varphi)
=
\frac{1}{8\,S(\varphi)}
\sum_{n=1}^{5}
\frac{(\varphi^{2n}+1)^3\,\bigl|\varphi^{2n}-1\bigr|}{\varphi^{4n}},
\]
with
\[
S(\varphi)
=
1+\frac18\sum_{n=1}^{5}\frac{(\varphi^{2n}+1)^4}{\varphi^{4n}},
\]
which is precisely \eqref{eq:Dob-Coff-Q(5)-a}.

For general $s\in\{1,2,3,4,5\}$ the matrices $\mathbb{P}^{(s)}(\varphi)$ and
$\mathbb{Q}^{(s)}(\varphi)$ have the same structure, with $10$ replaced by $2s$
and the sums run from $n=1$ to $s$. Repeating the above arguments with $5$
replaced by $s$ yields the general formulas
\eqref{eq:tauP-s-phi-def}--\eqref{eq:tauQ-s-phi-def}. The symmetry
\eqref{eq:tau-symmetry}, the bounds \eqref{eq:tau-bounds}, the values at
$\varphi=1$ in \eqref{eq:tau-at-1}, and the asymptotics
\eqref{eq:tau-asymptotics}--\eqref{eq:tauQ-asymptotics} follow immediately from
these explicit representations. This completes the proof.
\end{proof}

\end{document}
